\title{\textbf{基于多源预测与线性规划的微网协同调度}}
\date{}
\maketitle
\begin{abstract}
针对含光伏、储能和外部电网的小区微网购电调度问题，本文以10分钟为基本时段，
建立统一的电量流与储能状态模型，并依次研究确定性调度、供需不确定、日内预报更新和波动电价下的购电策略。

\textbf{针对问题一}，建立\textbf{全时段电量流线性规划模型}，统筹外网购电、光伏利用和储能充放电。
代表日购电费用为 35126.85元，较峰谷策略和光伏空间预留策略分别降低 \textbf{8.23\%} 和 \textbf{7.43\%}。

\textbf{针对问题二}，建立\textbf{周期预测—近期加权准备—跨日线性规划模型}。
分别预测负载和光伏，并根据近期净需求误差设置安全准备量，以降低5倍紧急购电风险。
2025年2月至12月总费用为 \textbf{1358.38万元}，较等权准备方案节省 44.66万元。

\textbf{针对问题三}，建立\textbf{历史融合预报—动态安全准备—日内滚动规划模型}，在0:00、6:00、12:00和18:00根据新预报和实际库存调整未执行订单。
全年总费用为 \textbf{1314.65万元}；相较初始融合方案节省 \textbf{94.22万元}，后发官方预报进一步带来 \textbf{21.78万元} 的费用改善。

\textbf{针对问题四}，将波动电价纳入预测与调度。
零点一次订电采用滚动提升树—线性规划模型，日内调整采用滚动岭回归—线性规划模型，两种方案总费用分别为 \textbf{1424.86万元} 和 \textbf{1382.29万元}。

模型检验表明，各方案均满足供需平衡及储能容量、功率约束。
本文形成了由\textbf{确定性调度—供需不确定性—日内信息更新—价格不确定性}逐步扩展的微网购电优化框架，可为类似微网的滚动购电与储能调度提供参考。
\end{abstract}
\textbf{关键词：} 微网调度；储能优化；线性规划；滚动优化；波动电价

\newpage